\section{The one-step carrier criterion}
\label{sec:one-step}

For $h\in H$, define its carrier height by
\[
 \operatorname{ht}(h):=
 \min\{\dim W:h\text{ occurs in }\mathcal B(W),\ W
             \text{ smooth projective}\}.
 \tag{4.1}\label{eq:height}
\]
Equivalently, $h\in H_{\le e}$ if and only if
\(\operatorname{ht}(h)\le e\).  This is the dimension of a smallest carrier,
not the dimension of a spectral-cover component or of the base of an atomic
$F$-bundle.

\begin{corollary}[Rank-two one-step criterion]
\label{cor:one-step}
Let $Z$ be a smooth projective $n$-fold with $n\ge1$, let $E\to Z$ be
a vector bundle of rank two, and let $h$ occur in \(\mathcal B(Z)\).  If
\[
 \operatorname{ht}(h)>n-1,
 \tag{4.2}\label{eq:height-bound}
\]
then \(\PP_Z(E)\) is irrational.  In particular, a top-height class
\(\operatorname{ht}(h)=n\) obstructs rationality after this one projective
step.
\end{corollary}

\begin{proof}
The total dimension is $d=n+1\ge2$.  By
\eqref{eq:ledger-projective},
\[
 \mathcal B(\PP_Z(E))=2\mathcal B(Z),
\]
so the coefficient of $h$ remains positive.  Condition
\eqref{eq:height-bound} says $h\notin H_{\le d-2}$.  Hence
\[
 \rho_{d-2}(\mathcal B(\PP_Z(E)))\ne0,
\]
and Theorem~\ref{thm:birational-quotient} applies.
\end{proof}

For Hodge atoms, Corollary~\ref{cor:one-step} is the rank-two specialization
of the standard carrier-filtration criterion.  The role of the categorical
construction is exact but modest: it proves that the chemical formula and
every lawful marker are folds of one occurrence ledger.  It does not prove
that a particular Hodge atom has the required carrier height; that remains a
geometric exclusion statement.

The relation with the direct-QDM specialization of \cite{Epilogue} can now be
stated without conflating their objects.
\begin{center}
\begin{tabular}{@{}>{\raggedright\arraybackslash}p{0.20\textwidth}
                    >{\raggedright\arraybackslash}p{0.34\textwidth}
                    >{\raggedright\arraybackslash}p{0.34\textwidth}@{}}
\toprule
Datum & Hodge-atom specialization & Direct-QDM specialization \\
\midrule
Local object & component of the reduced unramified Euler spectral cover &
  typed generic numerical small-QDM block \\
Generated relation & elementary geometric correspondences &
  regular block isomorphisms supplied by QDM comparison \\
Effective fold & chemical formula and dimension quotient &
  chosen numerical block marker \\
One-step input & carrier-height exclusion & low-dimensional marker nullity \\
\bottomrule
\end{tabular}
\end{center}
Both columns instantiate Theorem~\ref{thm:ledger}.  They do not use the same
groupoid, and a class in one column is not automatically a class in the
other.  This typed distinction is what permits the companion to clarify the
standard atom specialization without enlarging the proof of cubic one-step
irrationality.

The general projective-bundle correspondence in
Definition~\ref{def:provider} is part of the standard Hodge-atom quotient and
is also needed to evaluate projective space.  The stabilization consequence
of this note is only Corollary~\ref{cor:one-step}, for a rank-two bundle; no
claim is made for two or more projective stabilization factors.
